\chapter{The \texttt{workflow://} Schema: Data-Oriented Dependency Inference}
\label{ch:schema}

The \texttt{workflow://} schema is DAGonStar's mechanism for expressing
file-level dependencies between tasks without requiring the workflow author
to state them explicitly. This chapter gives the schema a complete, formal
treatment, because it is the single mechanism that later unifies file-based
scientific software (\Cref{ch:external-software}), LLM tasks
(\Cref{ch:llm-tasks}), Web tasks (\Cref{ch:web-tasks}), and Native Python
tasks (\Cref{ch:native-tasks}) under one dependency model.

\section{Purpose}

The schema lets a command refer to a file produced by another task. DAGonStar
infers the dependency from that reference and stages the file before the
consumer task runs --- turning a textual reference inside a shell command (or
JSON request, or function-argument mapping) into a structural edge of the
workflow DAG.

\section{Formal Shape}

\begin{definitionbox}{\texttt{workflow://} URI Grammar}
\begin{lstlisting}[basicstyle=\ttfamily\small, frame=none]
workflow://<workflow-name>/<task-name>/<path-inside-task-working-directory>
\end{lstlisting}
For references inside the \emph{same} workflow, \texttt{<workflow-name>} may
be omitted, yielding the common shorthand:
\begin{lstlisting}[basicstyle=\ttfamily\small, frame=none]
workflow:///<task-name>/<path-inside-task-working-directory>
\end{lstlisting}
\end{definitionbox}

\begin{lstlisting}[caption={Representative \texttt{workflow://} references.}]
workflow:///A/output/f1.txt
workflow://Forecast/preprocess/out/forcing.nc
workflow://Observations/download/data/profile.csv
\end{lstlisting}

\subsection{Components}

\paragraph{Scheme.} The literal scheme \texttt{workflow://}, defined in the
implementation as the class constant \texttt{Workflow.SCHEMA = "workflow://"}.

\paragraph{Workflow name.} Identifies the \emph{producer} workflow. Optional
for references inside the current workflow (\texttt{workflow:///producer/output.dat}),
mandatory for cross-workflow references
(\texttt{workflow://OtherWorkflow/producer/output.dat}; see
\Cref{ch:meta-async} for meta-workflow composition).

\paragraph{Task name.} The producer task's \texttt{name} attribute, e.g.\
\texttt{prepare} in \texttt{workflow:///prepare/input.nc}.

\paragraph{Path.} The remaining path segment, interpreted \emph{relative to
the producer task's working directory}. For example,
\texttt{workflow:///prepare/output/result.nc} resolves to
\texttt{<prepare working directory>/output/result.nc}.

\section{Dependency Inference Algorithm}

When \texttt{Workflow.make\_dependencies()} is called, every task runs its
\texttt{pre\_run()} method. \texttt{pre\_run()} scans the task's command
string for \texttt{workflow://} occurrences, and, for each occurrence:

\begin{enumerate}[leftmargin=1.6em]
  \item parses the workflow name, task name, and file path;
  \item finds the producer task (locally, or via the optional DAGon service
    for cross-workflow references when no local match exists);
  \item appends the consumer task to the producer's \texttt{nexts};
  \item appends the producer task to the consumer's \texttt{prevs};
  \item increments the producer's \texttt{reference\_count}.
\end{enumerate}

Consequently, the single-line command

\begin{lstlisting}
DagonTask(TaskType.BATCH, "B", "cat workflow:///A/output/f1.txt")
\end{lstlisting}

deterministically creates the graph edge $A \to B$, with no separate
\texttt{add\_dependency\_to()} call required.

\begin{figure}[h]
\centering
\begin{tikzpicture}[
  node distance=14mm and 16mm,
  task/.style={rectangle, rounded corners, draw=dagonblue, very thick,
    fill=blue!4, minimum width=26mm, minimum height=10mm, font=\sffamily\small,
    align=center, text width=30mm},
  note/.style={font=\sffamily\scriptsize, text width=44mm, align=left},
  arr/.style={-{Stealth[length=2.6mm]}, thick, dagonblue}
]
  \node[task] (A) {Task A\\ \texttt{output/f1.txt}};
  \node[task, right=of A] (B) {Task B\\ \texttt{cat workflow:///A/...}};
  \draw[arr] (A) -- node[above,font=\sffamily\scriptsize]{parsed by \texttt{pre\_run()}} (B);
  \node[note, below=2mm of A, xshift=8mm] {1. \texttt{make\_dependencies()} scans B's command\\
    2. finds reference to A\\
    3. records $A \to B$ in \texttt{prevs}/\texttt{nexts}\\
    4. stages \texttt{f1.txt} into B's \texttt{.dagon/inputs/}\\
    5. rewrites B's command to use the staged path};
\end{tikzpicture}
\caption{Dependency inference and command rewriting triggered by a single
\texttt{workflow://} reference.}
\label{fig:schema-resolution}
\end{figure}

\section{Command Rewriting}

Before execution, \texttt{Task.pre\_process\_command()} stages the referenced
file and rewrites the command body: the producer reference is replaced by a
path inside the consumer task's working directory, under

\begin{lstlisting}[basicstyle=\ttfamily\small, frame=none]
.dagon/inputs/<workflow-name>/<task-name>/...
\end{lstlisting}

This lets the external program --- which may be an arbitrary, unmodified
scientific executable --- read what it sees as an ordinary, local filesystem
path, with no awareness that the file originated in a sibling task's scratch
directory.

\section{Whitespace Boundary Rule}
\label{sec:whitespace-rule}

\begin{cautionbox}
The current parser scans from \texttt{workflow://} until the next whitespace
character or the end of the command string. Paths containing spaces are
therefore \emph{not} safely supported in \texttt{workflow://} references.

\textbf{Recommended:} \texttt{workflow:///A/output/result.nc}

\textbf{Avoid:} \texttt{workflow:///A/output/result with spaces.nc}
\end{cautionbox}

This constraint should be treated as a hard naming rule for any scientific
data product that may ever be referenced through the schema: avoid spaces in
output file and directory names produced by any task.

\section{Same-Workflow and Cross-Workflow Examples}

\begin{lstlisting}[caption={Same-workflow producer/consumer pair.}]
task_a = DagonTask(TaskType.BATCH, "A", "mkdir -p output; echo data > output/f1.txt")
task_b = DagonTask(TaskType.BATCH, "B", "cat workflow:///A/output/f1.txt > copy.txt")
\end{lstlisting}

\begin{lstlisting}[caption={Cross-workflow reference, used by meta-workflows.}]
task = DagonTask(
    TaskType.BATCH,
    "run_model",
    "cat workflow://Observations/collect/out/forcing.txt > model_input.txt",
)
\end{lstlisting}

Cross-workflow behavior underlies the meta-workflow and transversal patterns
developed in \Cref{ch:meta-async}, where it may involve \texttt{DAG\_TPS} or
the optional DAGon service.

\section{Interaction with Data Staging}

The schema identifies \emph{what} must be staged; \emph{how} it is staged is
controlled by the workflow's or task's data mover
(\Cref{sec:data-staging-architecture}): \texttt{COPY}, \texttt{LINK},
\texttt{SCP}, \texttt{GRIDFTP}, or \texttt{SKYCDS}. For local workflows,
\texttt{COPY} is the safest default; specialized movers require external
services or remote configuration.

\section{Extension to Non-File-Based Tasks}

A defining property of the engine's later AI--HPC integration is that the
\texttt{workflow://} schema is \emph{not} restricted to shell commands.
\texttt{TaskType.LLM} parses \texttt{workflow://} values in both its JSON
prompt and its \texttt{input\_files} mapping, infers the producer dependency,
and copies local UTF-8 text inputs into its \texttt{.dagon/inputs} folder
before including their content in the API request body
(\Cref{ch:llm-tasks}). \texttt{TaskType.WEB} similarly parses recursive
\texttt{workflow://} values inside an HTTP request specification
(\Cref{ch:web-tasks}), and \texttt{TaskType.NATIVE} accepts them as Python
function-argument bindings (\Cref{ch:native-tasks}). The schema is, in this
precise sense, the engine's \emph{universal} dependency language.

\section{Best Practices}

\begin{itemize}[leftmargin=1.4em]
  \item Use stable task names; the schema and checkpointing both key on them.
  \item Use simple path names without spaces (\Cref{sec:whitespace-rule}).
  \item Write producer outputs under explicit, named directories such as
    \texttt{output/}.
  \item Reference exact files rather than broad directory assumptions where
    possible.
  \item Always call \texttt{workflow.make\_dependencies()} before
    \texttt{workflow.run()} when the dependency graph should be inspected or
    validated ahead of execution --- although, as \Cref{ch:meta-async} notes,
    \texttt{run()} and \texttt{launch()} build it automatically if it has
    not already been built.
  \item Keep cross-workflow references rare and well documented.
\end{itemize}

\section{Common Errors}

\begin{description}[leftmargin=1.8em,style=nextline]
  \item[Missing producer task] If a referenced task is not found and the
    DAGon service is disabled, DAGonStar raises a connection/service-related
    error, because it cannot resolve the reference externally.
  \item[Dependency cycle] If \texttt{workflow://} references create a cycle,
    graph validation fails (\Cref{sec:dag-theory}).
  \item[Wrong path inside producer task] The dependency may be structurally
    valid, yet staging or downstream execution fails if the producer did not
    actually create the referenced file --- a reminder that the schema
    expresses \emph{intended} data lineage, not a guarantee of file
    existence.
  \item[Path contains spaces] The parser uses whitespace as a token
    boundary; see \Cref{sec:whitespace-rule}.
\end{description}

\section{The Schema as a Naming Scheme: A Formal View}
\label{sec:schema-formal}

It is useful to read the \texttt{workflow://} grammar of
\Cref{ch:schema} through the lens of \emph{naming schemes} in distributed
systems, a perspective that clarifies precisely what guarantee the schema
does, and does not, provide.

A naming scheme maps names --- here, \texttt{workflow://} URIs --- to
resources --- here, files inside task scratch directories. Two properties
are commonly demanded of a useful naming scheme: \emph{location
transparency}, meaning a name does not encode the physical location of the
resource it refers to, only its logical identity; and \emph{resolvability},
meaning the system can deterministically map a name to a current physical
location. The \texttt{workflow://} schema satisfies location transparency by
construction: a reference such as \texttt{workflow:///prepare/output/result.nc}
says nothing about \emph{where} \texttt{prepare}'s scratch directory
physically resides --- under \texttt{/tmp/}, under a configured alternate
\texttt{scratch\_dir\_base} (\Cref{ch:configuration}), or, for a remote task,
on a different host entirely. The mapping from this logical name to a
physical path is resolved at \texttt{pre\_process\_command()} time
(\Cref{sec:execution-lifecycle}), using only information already present in
the in-memory task graph: the producer task's recorded \texttt{working\_dir}.

This formal framing also clarifies why the schema's resolvability is
\emph{intra-workflow-execution} only: a \texttt{workflow://} reference is
meaningful only in the context of a graph that DAGonStar itself constructed
and whose tasks' working directories it itself tracked. A reference copied
out of one workflow run's source code and pasted, unmodified, into an
unrelated script with no corresponding producer task would not resolve to
anything --- the name carries no inherent physical meaning outside the
dependency graph that gives it one. This is, again, exactly the
\emph{Missing producer task} failure mode already catalogued in
\Cref{ch:schema}'s ``Common Errors'' section, now understood as a structural
consequence of how the naming scheme is resolved rather than as an
arbitrary implementation limitation.

\section{Comparison with Alternative Dependency-Expression Mechanisms}

Having developed the schema's mechanics in full, it is worth comparing it
explicitly against the two dependency-expression alternatives a workflow
engine could instead have chosen, to clarify what specific trade-off
DAGonStar's designers accepted.

\textbf{Configuration-file dependency declaration} --- in which a separate
manifest (a YAML or JSON file, say) lists, for every task, the names of its
predecessors --- decouples the dependency declaration from the command
string entirely. This has the advantage of being inspectable by tooling that
does not parse arbitrary shell syntax, but it reintroduces precisely the
synchronization risk \Cref{sec:task-vs-data} identified for explicit,
asserted dependencies: a manifest entry and the command that actually
consumes a file can silently drift apart, because nothing forces the two to
agree. \texttt{workflow://} dependency inference avoids this drift
structurally, by deriving the dependency from the very same string the
consumer's command will execute.

\textbf{Implicit, timestamp-based dependency inference} --- the strategy
classical build systems such as \texttt{make} use, inferring that a target
is stale if any prerequisite file's modification time is newer --- requires
no explicit reference syntax at all, but it depends on a shared, reliable
filesystem clock and breaks down precisely in the heterogeneous,
multi-host, multi-backend setting (\Cref{ch:backends}) DAGonStar is designed
for: a remote task's output file's modification timestamp, as observed
through SSH, may not be directly comparable to a local task's, especially
across hosts with imperfectly synchronized clocks. \texttt{workflow://}
inference sidesteps the clock-synchronization problem entirely, because it
never inspects timestamps; the dependency is structural, derived once at
\texttt{make\_dependencies()} time, and the engine's own
\texttt{Status} state machine (\Cref{sec:status-model}), not the filesystem
clock, is the sole authority on whether a producer has actually finished.

\begin{practicebox}
The schema's chosen strategy --- embedding the dependency reference directly
inside the artifact (command, prompt, or request specification) that
consumes it --- trades a small amount of syntactic novelty (a workflow
author must learn the \texttt{workflow://} grammar) for the elimination of
an entire class of synchronization bugs that both configuration-file and
timestamp-based alternatives are structurally prone to. This trade-off is
visible again, in a different guise, every time \Cref{ch:llm-tasks,ch:web-tasks,ch:native-tasks}
show the same embedded-reference strategy extended into a JSON prompt, an
HTTP request specification, and a Python argument-binding dictionary,
respectively.
\end{practicebox}

\section{Provenance Beyond the Single Reference}

A single \texttt{workflow://} reference establishes one edge of a
provenance graph; it is worth being explicit about what larger structure
emerges once a realistic workflow contains many such references, because
this larger structure is the actual scientific deliverable a DAGonStar
workflow produces, alongside its numerical results.

Define the \emph{provenance graph} of a completed workflow run as the
directed graph whose nodes are the workflow's tasks, annotated with their
recorded \texttt{working\_dir}, generated \texttt{.dagon/launcher.sh}
(\Cref{ch:lessons-03-04}), and final \texttt{Status}, and whose edges are
exactly the dependencies \texttt{make\_dependencies()} constructed --- the
same structure already developed as the dependency DAG, but now read
\emph{retrospectively}, after execution, as a historical record rather than
a forward scheduling plan. Because every edge in this graph corresponds to
an actual, staged file copy (under the chosen \texttt{DataMover}) and a
rewritten command, the provenance graph is, in a precise sense, more
complete than the kind of provenance a researcher typically reconstructs by
hand from a shell history or a lab notebook: it records not just \emph{that}
one step followed another, but the \emph{exact} file each step consumed,
the exact path it was staged to, and the exact, fully-resolved command that
ran.

\begin{example}
Consider a three-task chain \texttt{download} $\to$ \texttt{preprocess}
$\to$ \texttt{simulate}, of the kind \Cref{ch:lessons-05-07} discusses in
the context of checkpointing. After a successful run, a researcher who
receives only the three scratch directories --- without the original Python
source --- can nonetheless reconstruct the full provenance chain by reading
each task's \texttt{.dagon/launcher.sh} in turn: \texttt{simulate}'s
launcher will reference a path under
\texttt{.dagon/inputs/<workflow>/preprocess/...}, which can be traced back
to \texttt{preprocess}'s own output directory, and so on. This
reconstructibility-from-artifacts-alone property is a direct, mechanical
consequence of the command-rewriting behavior described in \Cref{ch:schema},
and is part of why this book treats the generated launcher script, first
introduced in \Cref{ch:lessons-03-04}, as a first-class scientific
artifact in its own right, not merely an implementation detail.
\end{example}

\section{Chapter Summary}

The \texttt{workflow://} schema converts a textual file reference embedded in
a task's command, prompt, request specification, or argument binding into a
structural DAG edge, a staged local file, and a rewritten, self-contained
command. It is the single dependency language shared by every execution
backend covered in this book, from a local \texttt{BATCH} task in
\Cref{ch:lessons-01-02} to an \texttt{LLM} task in \Cref{ch:llm-tasks}. This
chapter additionally framed the schema formally as a location-transparent
naming scheme, compared it against configuration-file and timestamp-based
dependency-inference alternatives, and showed how a workflow's accumulated
\texttt{workflow://} references constitute a reconstructible provenance
graph that survives independently of the original source code.

\begin{exercise}
Using the provenance-graph notion of this chapter's final section, describe
precisely what information would be \emph{lost}, and what would still be
recoverable, if a workflow author replaced every \texttt{workflow://}
reference in a five-task pipeline with explicit \texttt{add\_dependency\_to()}
calls plus hard-coded absolute paths copied by hand from a previous run's
scratch directories. Relate your answer to \Cref{sec:task-vs-data}.
\end{exercise}
